\section{Síntesis y ampliación}

Esta clase parte de una pregunta sencilla pero incisiva: ¿cómo puede una red neuronal de conexiones fijas representar una part-whole hierarchy que es distinta para cada entrada? La respuesta de GLOM no asigna node hardware de forma dinámica. En su lugar, deja que la activity de la position-by-level grid forme islands mediante recurrent interaction. Los vectores aproximadamente coincidentes de una misma island representan un parse node. Las learned functions entre el nivel superior y el inferior expresan la part-whole prediction, la lateral attention expresa la agreement dentro del mismo nivel y el location-conditioned neural field hace que un mismo whole genere parts distintas en posiciones distintas.

\subsection{La clase completa condensada en seis pasos causales}

\begin{enumerate}
  \item \textbf{Dynamic state replaces dynamic memory allocation}: la estructura específica de cada entrada se escribe en la activity, no modificando los weights;
  \item \textbf{Coordinate-aware prediction proposes hierarchy}: las parts votan hacia arriba por el whole y el whole predice hacia abajo las parts;
  \item \textbf{Local similarity attention forms islands}: las same-level hypotheses que ya son similares se refuerzan entre sí;
  \item \textbf{Recurrent settling delays commitment}: la inference de varios pasos admite ambiguity, corrección y top-down feedback;
  \item \textbf{Reconstruction and consensus train the loop}: el prediction error externo y la co-distillation interna moldean conjuntamente la representación;
  \item \textbf{Location-conditioned decoding preserves weight sharing}: un mismo whole vector produce salidas distintas para cada coordinate.
\end{enumerate}

\begin{importantbox}{Conclusión didáctica final}
Lo que más vale conservar de GLOM no es una update rule aún no validada, sino una representation discipline: si afirmas que un modelo entiende la estructura de los objetos, debes explicar dónde existen los node, cómo se codifica la part-whole relation, cómo se preserva la ambiguity, cómo se forma el binding, cómo se entrena la estructura mediante la loss y qué intervention puede demostrar que el cómputo posterior realmente utiliza esa estructura.
\end{importantbox}

\subsection{Seis errores frecuentes}

\begin{warningbox}{No presentes una propuesta de diseño como una historia de éxito}
\begin{enumerate}
  \item No digas que GLOM ya resolvió la object-centric vision; esta clase no presenta un benchmark completo;
  \item No tomes las flechas bidimensionales por la embedding dimension real ni por una pose matrix explícita;
  \item No equipares directamente una island con una ground-truth segmentation mask;
  \item No equipares consensus con truth: la agreement también puede ser un shared error;
  \item No tomes la Hough analogy como una complexity advantage demostrada;
  \item No equipares el location-conditioned decoder con un NeRF/modern world model completo.
\end{enumerate}
\end{warningbox}

\subsection{Una conexión prudente con los sistemas actuales}

Las object-centric representation, los slot-based model, el iterative refinement, las equivariant network, las implicit neural representation, los video world model y la sparse hierarchical attention actuales abordan, cada uno, problemas parciales de GLOM. La comparación verdaderamente valiosa no se fija en los nombres, sino que revisa punto por punto: si hay un dynamic entity state, si se trata la pose de forma explícita, si se conservan multiple hypotheses, si se admite top-down correction, si existe una part-whole relation legible y si la estructura se ha validado mediante causal intervention.

\teachervoice{Antes de terminar, Hinton dice que fue una talk compleja y que quizá su mejor uso sea animar a todos a leer el artículo largo. Este cierre hace eco del «vaporware» de la apertura: no anuncia que el problema esté resuelto, sino que invita a los lectores a convertir una representation idea en un modelo y en experimentos más rigurosos.}

\subsection{Lecturas adicionales}

\begin{center}
\footnotesize
\begin{tabularx}{\textwidth}{>{\raggedright\arraybackslash}X>{\raggedright\arraybackslash}X}
\href{https://arxiv.org/abs/2102.12627}{Hinton, \emph{GLOM}}: documento de diseño completo. &
\href{https://arxiv.org/abs/1710.09829}{Sabour et al., \emph{Capsules}}: antecedentes de routing. \\[0.25em]
\href{https://arxiv.org/abs/2002.05709}{Chen et al., \emph{SimCLR}}: contrastive scaffold. &
\href{https://arxiv.org/abs/1706.03762}{Vaswani et al., \emph{Transformer}}: línea base de attention. \\[0.25em]
\href{https://arxiv.org/abs/1810.04805}{Devlin et al., \emph{BERT}}: masked reconstruction. &
\href{https://arxiv.org/abs/1804.03235}{Anil et al., \emph{Online Distillation}}: co-distillation. \\[0.25em]
\href{https://arxiv.org/abs/2003.08934}{Mildenhall et al., \emph{NeRF}}: coordinate-conditioned field. & \\
\end{tabularx}
\end{center}

\end{document}
